\subsection{Two-dimensional results: conclusions}
\label{sec:results-2d-conclusions}

The three nonlinear strategies give almost identical convergence orders for a
fixed spatial, temporal, and mass-matrix configuration. Therefore, in this MMS
case the main distinction between \code{Picard}, \code{diagonalIion}, and
\code{JFNK} is computational cost rather than asymptotic accuracy. The clearest
spatial orders are obtained with \code{crankNicolson}, \code{consistent} mass,
and $\Delta t=10^{-3}$; for $\Delta t=10^{-2}$ the finest-mesh slopes of the
high-order state diagnostics are still visibly affected by time error.

\paragraph{Picard.}
\code{Picard} is the most direct baseline and reproduces the expected
high-order trends when the rest of the discretisation is chosen appropriately.
For example, with \code{p3} voltage reconstruction, \code{p2} source/state
quadrature, \code{crankNicolson}, \code{consistent} mass, and
$\Delta t=10^{-3}$, the fitted orders are
$3.538$ for the cell-centred $\Vm$, $3.992$ for $\Vm^G$, and about $4.0$ for
the reconstructed state errors $u_1^G$ and $u_2^G$
(Table~\ref{tab:results-2d-picard-1p0em03}). Increasing the source/state level
from \code{p2} to \code{p3} changes these orders negligibly but increases the
runtime. The disadvantage is efficiency: the same \code{p3}/\code{p2},
\code{crankNicolson}, \code{consistent}, $\Delta t=10^{-3}$ sweep takes
$1730.1$~s with \code{Picard}, compared with $798.27$~s with \code{JFNK}.

\paragraph{\code{diagonalIion}.}
\code{diagonalIion} produces essentially the same orders as \code{Picard}. Its
diagonal treatment of the local ionic-current sensitivity can help the
nonlinear voltage equation behave more implicitly, but for the present
manufactured problem this does not translate into a systematic gain in either
accuracy or runtime. In the same \code{p3}/\code{p2}, \code{crankNicolson},
\code{consistent}, $\Delta t=10^{-3}$ case, the orders again are $3.538$ for
$\Vm$, $3.992$ for $\Vm^G$, and approximately $4.0$ for $u_1^G,u_2^G$, while
the accumulated runtime is $1661.8$~s
(Table~\ref{tab:results-2d-diagonaliion-1p0em03}). Thus, for these 2D MMS
runs, \code{diagonalIion} is mainly a useful robustness option, not the best
cost--accuracy choice.

\paragraph{JFNK.}
\code{JFNK} gives the same convergence orders as the other methods but is the
most efficient nonlinear coupling in these results. For the recommended
\code{p3}/\code{p2}, \code{crankNicolson}, \code{consistent},
$\Delta t=10^{-3}$ configuration, it obtains $3.539$ for $\Vm$, $3.993$ for
$\Vm^G$, $4.003$ for $u_1^G$, and $4.010$ for $u_2^G$ over $N=10$--$100$; over
$N=70$--$100$, the reconstructed diagnostics remain close to fourth order
(Table~\ref{tab:results-2d-jfnk-1p0em03}). The accumulated runtime is
$798.27$~s, less than half the corresponding \code{Picard} and
\code{diagonalIion} runs, while the summed peak-RSS cost is comparable. This
makes \code{JFNK} the best nonlinear method for the 2D manufactured study.

\paragraph{Temporal scheme.}
For high-order spatial verification, the second-order \code{crankNicolson}
scheme is preferable. The first-order time error of \code{backwardEuler} masks
the high-order spatial behaviour on the fine meshes; for example, the
\code{p3}/\code{p2}, \code{consistent}, $\Delta t=10^{-3}$ \code{backwardEuler}
case has good full-range reconstructed slopes but the $N=70$--$100$ slopes of
$\Vm^G$, $u_1^G$, and $u_2^G$ collapse to values below one in
Table~\ref{tab:results-2d-jfnk-1p0em03}. With \code{crankNicolson}, the same
spatial configuration recovers fourth-order reconstructed behaviour.
\code{backwardEuler} is useful as a robustness or comparison run, but it should
not be used as the primary evidence for high-order spatial convergence.

\paragraph{Mass matrix.}
The \code{consistent} mass matrix is the better choice for the high-order
configurations. The effect is especially clear for \code{p3}: with
\code{JFNK}, \code{crankNicolson}, $\Delta t=10^{-3}$, and \code{p3}/\code{p2},
the \code{consistent} mass case gives full-range orders $3.539$, $3.993$,
$4.003$, and $4.010$ for $\Vm$, $\Vm^G$, $u_1^G$, and $u_2^G$, whereas the
\code{lumped} case gives $2.369$, $3.091$, $3.441$, and $3.435$
(Table~\ref{tab:results-2d-jfnk-1p0em03}). The lumped matrix remains adequate
for low-order or second-order comparisons, but it degrades the high-order
\code{p3} regime.

\paragraph{Matching the voltage and source/state orders.}
The source/state quadrature should usually be matched to the voltage
reconstruction level, but the data suggest one important practical refinement:
\begin{itemize}[leftmargin=*]
\item \code{NO} voltage should be paired with \code{NO} source/states. Higher
source quadrature does not improve the second-order voltage result and only
adds cost.
\item \code{p1} voltage should be paired with at least \code{p1}
source/states. The \code{p1}/\code{NO} cases under-resolve the state coupling,
whereas \code{p1}/\code{p1} gives the expected second-order reconstructed
state behaviour.
\item \code{p2} voltage should be paired with \code{p2} source/states.
Although the formal interpolation expectation for \code{p2} is third order, the
cell-centred finite-volume unknown often converges with order near two or
slightly above two; in these data the reconstructed Gauss-point diagnostics and
source/state quantities show the clearer benefit of the quadratic
reconstruction.
\item \code{p3} voltage should be paired with \code{p2} source/states for the
best cost--accuracy balance. The \code{p3}/\code{p3} option is more symmetric
and remains a conservative verification choice, but in the present 2D results
it gives essentially the same orders as \code{p3}/\code{p2} at higher runtime.
\end{itemize}

\paragraph{Recommended 2D configuration.}
Considering both convergence and efficiency, the best overall configuration in
the current 2D MMS data is
\[
  \boxed{
  \code{JFNK}
  + \code{crankNicolson}
  + \code{consistent mass}
  + \Delta t=10^{-3}
  + \Vm:\code{p3}
  + \Iion/\mathbf{u}:\code{p2}}
\]
This choice reaches fourth-order behaviour in the reconstructed voltage and
state diagnostics, keeps the cell-centred voltage close to fourth order on the
fine meshes, and is substantially faster than the equivalent \code{Picard} and
\code{diagonalIion} runs. For a cheaper second-order study, the consistent
recommendation is \code{JFNK} with \code{crankNicolson} and matched
\code{NO}/\code{NO} or \code{p1}/\code{p1} reconstruction, depending on whether
only the standard finite-volume method or the first high-order reconstruction
level is being tested.
